%%%%%%%%%%%%%%%%
%%% exod 33 %%%
%%%%%%%%%%%%%%%%

\Chapter{33}

\Verse{1}
{
	\hJ{וַיְדַבֵּר יְהֹוָה אֶל מֹשֶׁה לֵךְ עֲלֵה מִזֶּה אַתָּה וְהָעָם אֲשֶׁר הֶעֱלִיתָ מֵאֶרֶץ מִצְרָיִם אֶל הָאָרֶץ אֲשֶׁר נִשְׁבַּעְתִּי לְאַבְרָהָם לְיִצְחָק וּלְיַעֲקֹב לֵאמֹר לְזַרְעֲךָ אֶתְּנֶנָּה׃}
}
{
	\eJ{
		33 And YHWH spoke to Moses: “Go. Go up from here, you and
		the people whom you brought up from the land of Egypt, to
		the land that I swore to Abraham, to Isaac, and to Jacob,
		saying, ‘I’ll give it to your seed.’
	}
}
{
}

\Verse{2}
{
	\hJ{וְשָׁלַחְתִּי לְפָנֶיךָ מַלְאָךְ וְגֵרַשְׁתִּי אֶת הַכְּנַעֲנִי הָאֱמֹרִי וְהַחִתִּי וְהַפְּרִזִּי הַחִוִּי וְהַיְבוּסִי׃}
}
{
	\eJ{
		And I’ll send an angel ahead of you—and I’ll drive out the
		Canaanite, the Amorite, and the Hittite and the Perizzite
		and the Hivite and the Jebusite—
	}
}
{
}

\Verse{3}
{
	\hJ{אֶל אֶרֶץ זָבַת חָלָב וּדְבָשׁ כִּי לֹא אֶעֱלֶה בְּקִרְבְּךָ כִּי עַם קְשֵׁה עֹרֶף אַתָּה פֶּן אֲכֶלְךָ בַּדָּרֶךְ׃}
}
{
	\eJ{
		to a land flowing with milk and honey, because I won’t go
		up among you, because you’re a hard-necked people, or else
		I’ll finish you on the way.”
	}
}
{
}

\Verse{4}
{
	\hJ{וַיִּשְׁמַע הָעָם אֶת הַדָּבָר הָרָע הַזֶּה וַיִּתְאַבָּלוּ וְלֹא שָׁתוּ אִישׁ עֶדְיוֹ עָלָיו׃}
}
{
	\eJ{
		And the people heard this bad thing, and they mourned, and
		not a man put his jewelry on him.
	}
}
{
}

\Verse{5}
{
	\hJ{וַיֹּאמֶר יְהֹוָה אֶל מֹשֶׁה אֱמֹר אֶל בְּנֵי יִשְׂרָאֵל אַתֶּם עַם קְשֵׁה עֹרֶף רֶגַע אֶחָד אֶעֱלֶה בְקִרְבְּךָ וְכִלִּיתִיךָ וְעַתָּה הוֹרֵד עֶדְיְךָ מֵעָלֶיךָ וְאֵדְעָה מָה אֶעֱשֶׂה לָּךְ׃}
}
{
	\eJ{
		And YHWH said to Moses, “Say to the children of Israel:
		‘You’re a hard-necked people. If I would go up among you
		for one moment, I would finish you! And now, put your
		jewelry down from on you, so I may know what I’ll do with
		you.’”
	}
}
{
%	\footnote{
%		put your jewelry down. Did they not already take it off (v. 4)?! But
%		here the text appears to refer to permanent removal. Thus in v. 6 we
%		learn that they “divested their jewelry from Mount Horeb”—meaning: from
%		Mount Horeb on.
%	}
}

\Verse{6}
{
	\hJ{וַיִּתְנַצְּלוּ בְנֵי יִשְׂרָאֵל אֶת עֶדְיָם מֵהַר חוֹרֵב׃}
}
{
	\eJ{
		And the children of Israel divested their jewelry from
		Mount Horeb.
	}
}
{
}

\Verse{7}
{
	\hE{וּמֹשֶׁה יִקַּח אֶת הָאֹהֶל וְנָטָה לוֹ מִחוּץ לַמַּחֲנֶה הַרְחֵק מִן הַמַּחֲנֶה וְקָרָא לוֹ אֹהֶל מוֹעֵד וְהָיָה כּל מְבַקֵּשׁ יְהֹוָה יֵצֵא אֶל אֹהֶל מוֹעֵד אֲשֶׁר מִחוּץ לַמַּחֲנֶה׃}
}
{
	\eE{
		And Moses would take the tent and pitch it outside of the
		camp, going far from the camp, and he called it the Tent
		of Meeting. And it would be: everyone seeking YHWH would
		go out to the Tent of Meeting, which was outside of the
		camp.
	}
}
{
%	\footnote{
%		the Tent of Meeting This is the first mention of the tent where God will
%		communicate with Israel. It will be the central place of worship, the
%		only place where sacrifice is allowed, for all the years that they will
%		spend in the wilderness and for centuries after that in the promised
%		land, until its disappearance at the time of the destruction of the
%		Temple of Solomon. The problem is that the Tent of Meeting is not yet
%		built at this point. It is not set up until Exodus 40. Rashi and Ibn
%		Ezra therefore identified the tent in this verse as Moses’ own tent
%		(which has been mentioned already in 18:7). Ibn Ezra also says that
%		there are those who say that it is ’hel hammiškn, the Tent of the
%		Tabernacle. From the point of view of critical scholarship, the problem
%		is the result of the combining of two of the sources of the Torah: the
%		report of the moving and naming of the Tent is identified as coming from
%		E; the later report of the setting up of the Tent is identified as
%		coming from P. That is an explanation of how the text was composed, but
%		in understanding the text as it stands, I believe that Rashi’s and Ibn
%		Ezra’s explanation fits the text best. It appears that Moses’ tent
%		serves as the Tent of Meeting until the Tabernacle is set up. The
%		Tabernacle then replaces Moses’ tent as the Tent of Meeting. This is
%		significant because the transition of the sacred place of revelation
%		away from Moses’ personal residence reduces the impression of Moses’ own
%		proprietary role. There is a continuing issue of the people’s focusing
%		on Moses rather than on the deity, and the use of Moses’ residence as
%		the sacred place would contribute to their impression.
%	}
}

\Verse{8}
{
	\hE{וְהָיָה כְּצֵאת מֹשֶׁה אֶל הָאֹהֶל יָקוּמוּ כּל הָעָם וְנִצְּבוּ אִישׁ פֶּתַח אהֳלוֹ וְהִבִּיטוּ אַחֲרֵי מֹשֶׁה עַד בֹּאוֹ הָאֹהֱלָה׃}
}
{
	\eE{
		And it would be, when Moses would go out to the Tent, all
		the people would get up, and they would stand up, each one
		at the entrance of his tent, and they would look after
		Moses until he came to the Tent.
	}
}
{
}

\Verse{9}
{
	\hE{וְהָיָה כְּבֹא מֹשֶׁה הָאֹהֱלָה יֵרֵד עַמּוּד הֶעָנָן וְעָמַד פֶּתַח הָאֹהֶל וְדִבֶּר עִם מֹשֶׁה׃}
}
{
	\eE{
		And it would be, when Moses came to the Tent, the column
		of cloud would come down, and it would stand at the
		entrance of the Tent, and He would speak with Moses.
	}
}
{
}

\Verse{10}
{
	\hE{וְרָאָה כל הָעָם אֶת עַמּוּד הֶעָנָן עֹמֵד פֶּתַח הָאֹהֶל וְקָם כּל הָעָם וְהִשְׁתַּחֲווּ אִישׁ פֶּתַח אהֳלוֹ׃}
}
{
	\eE{
		And all the people would see the column of cloud standing
		at the entrance of the Tent, and all the people would get
		up and bow, each at the entrance of his tent.
	}
}
{
}

\Verse{11}
{
	\hE{וְדִבֶּר יְהֹוָה אֶל מֹשֶׁה פָּנִים אֶל פָּנִים כַּאֲשֶׁר יְדַבֵּר אִישׁ אֶל רֵעֵהוּ וְשָׁב אֶל הַמַּחֲנֶה וּמְשָׁרְתוֹ יְהוֹשֻׁעַ בִּן נוּן נַעַר לֹא יָמִישׁ מִתּוֹךְ הָאֹהֶל׃}
}
{
	\eE{
		And YHWH would speak to Moses face-to-face, the way a man
		speaks to his fellow man. And he would come back to the
		camp. And his attendant, Joshua, son of Nun, a young man,
		would not depart from inside the Tent.
	}
}
{
%	\footnote{
%		face-to-face. This does not mean that Moses has seen God’s face. He has
%		not seen it up to this point, and a few verses later God will tell him
%		explicitly that he will see God from behind, and “my face won’t be seen”
%		(33:23). For further evidence that this expression does not mean
%		literally seeing a face, see the comment on its occurrence in Deut
%		Footnote 5:4. Another occurrence of this expression in the Torah is when
%		Jacob says “I’ve seen God face-to-face” (Gen 32:31) after he has
%		struggled with “a man,” which refers to an angel (see the comments on
%		Genesis 32).
%	}
}

\Verse{12}
{
	\hJ{וַיֹּאמֶר מֹשֶׁה אֶל יְהֹוָה רְאֵה אַתָּה אֹמֵר אֵלַי הַעַל אֶת הָעָם הַזֶּה וְאַתָּה לֹא הוֹדַעְתַּנִי אֵת אֲשֶׁר תִּשְׁלַח עִמִּי וְאַתָּה אָמַרְתָּ יְדַעְתִּיךָ בְשֵׁם וְגַם מָצָאתָ חֵן בְּעֵינָי׃}
}
{
	\eJ{
		And Moses said to YHWH, “See, you say to me, ‘Bring this
		people up,’ and you haven’t made known to me whom you will
		send with me. And you’ve said, ‘I’ve known you by name,’
		and also ‘You’ve found favor in my eyes.’
	}
}
{
}

\Verse{13}
{
	\hJ{וְעַתָּה אִם נָא מָצָאתִי חֵן בְּעֵינֶיךָ הוֹדִעֵנִי נָא אֶת דְּרָכֶךָ וְאֵדָעֲךָ לְמַעַן אֶמְצָא חֵן בְּעֵינֶיךָ וּרְאֵה כִּי עַמְּךָ הַגּוֹי הַזֶּה׃}
}
{
	\eJ{
		And now, if I’ve found favor in your eyes, make your way
		known to me, so I may know you, so that I’ll find favor in
		your eyes. And see that this nation is your people.”
	}
}
{
%	\footnote{
%		so I may know you, so that I’ll find favor in your eyes. Moses’ words in
%		this passage seem unclear and even self-contradictory, but they are not.
%		Moses is asking God to tell him how He will help him lead the people to
%		the land. Moses quotes God as saying “I’ve known you” and “You’ve found
%		favor in my eyes,” but he indicates that he cannot really know God and
%		find favor with God unless God informs him what His way will be.
%		Footnote 33:13. And see that this nation is your people. Moses ends his
%		plea with an expression of caring for his people as well. In God’s words
%		(33:1) and in Moses’ quotation of those words (33:12), God speaks of
%		“the people,” but Moses continues the battle of the pronouns (see my
%		comment on 32:7) and pleads: “See that they’re your people.”
%	}
}

\Verse{14}
{
	\hJ{וַיֹּאמַר פָּנַי יֵלֵכוּ וַהֲנִחֹתִי לָךְ׃}
}
{
	\eJ{And He said, “My face will go, and I’ll let you rest.”}
}
{
}

\Verse{15}
{
	\hJ{וַיֹּאמֶר אֵלָיו אִם אֵין פָּנֶיךָ הֹלְכִים אַל תַּעֲלֵנוּ מִזֶּה׃}
}
{
	\eJ{
		And he said to Him, “If your face isn’t going, don’t take
		us up from here.
	}
}
{
}

\Verse{16}
{
	\hJ{וּבַמֶּה יִוָּדַע אֵפוֹא כִּי מָצָאתִי חֵן בְּעֵינֶיךָ אֲנִי וְעַמֶּךָ הֲלוֹא בְּלֶכְתְּךָ עִמָּנוּ וְנִפְלִינוּ אֲנִי וְעַמְּךָ מִכּל הָעָם אֲשֶׁר עַל פְּנֵי הָאֲדָמָה׃}
}
{
	\eJ{
		And by what then will it be known that I’ve found favor in
		your eyes, I and your people? Is it not by your going with
		us? And we’ll be distinguished, I and your people, from
		every people that is on the face of the earth.”
	}
}
{
%	\footnote{
%		distinguished. This refers to being different from the other peoples in
%		this one respect: that God’s presence is among them. It does not
%		necessarily imply distinction in the sense of being superior in any way.
%	}
}

\Verse{17}
{
	\hJ{וַיֹּאמֶר יְהֹוָה אֶל מֹשֶׁה גַּם אֶת הַדָּבָר הַזֶּה אֲשֶׁר דִּבַּרְתָּ אֶעֱשֶׂה כִּי מָצָאתָ חֵן בְּעֵינַי וָאֵדָעֲךָ בְּשֵׁם׃}
}
{
	\eJ{
		And YHWH said to Moses, “I’ll do this thing that you’ve
		spoken as well, because you’ve found favor in my eyes, and
		I’ve known you by name.”
	}
}
{
%	\footnote{
%		you’ve found favor in my eyes, and I’ve known you by name. God repeats
%		the exact phrases that Moses used as the basis of his case, thus
%		confirming His acceptance of Moses’ point.
%	}
}

\Verse{18}
{
	\hJ{וַיֹּאמַר הַרְאֵנִי נָא אֶת כְּבֹדֶךָ׃}
}
{
	\eJ{And he said, “Show me your glory!”}
}
{
%	\footnote{
%		Show me your glory. Using the Hebrew n’ particle of polite request,
%		Moses boldly dares to ask to be allowed to see God. And see how far
%		Moses has come: at the burning bush “Moses hid his face, because he was
%		afraid of looking at God” (Exod 3:6), but now, after a time of
%		acquaintance with God and after seeing and sometimes wielding divine
%		power, he has grown in intimacy with God and has grown in his own
%		confidence enough to ask God to see the divine form.
%	}
}

\Verse{19}
{
	\hJ{וַיֹּאמֶר אֲנִי אַעֲבִיר כּל טוּבִי עַל פָּנֶיךָ וְקָרָאתִי בְשֵׁם יְהֹוָה לְפָנֶיךָ וְחַנֹּתִי אֶת אֲשֶׁר אָחֹן וְרִחַמְתִּי אֶת אֲשֶׁר אֲרַחֵם׃}
}
{
	\eJ{
		And He said, “I shall have all my good pass in front of
		you, and I shall invoke the name YHWH in front of you. And
		I shall show grace to whomever I shall show grace, and I
		shall show mercy to whomever I shall show mercy.”
	}
}
{
}

\Verse{20}
{
	\hJ{וַיֹּאמֶר לֹא תוּכַל לִרְאֹת אֶת פָּנָי כִּי לֹא יִרְאַנִי הָאָדָם וָחָי׃}
}
{
	\eJ{
		And He said, “You won’t be able to see my face, because a
		human will not see me and live.”
	}
}
{
}

\Verse{21}
{
	\hJ{וַיֹּאמֶר יְהֹוָה הִנֵּה מָקוֹם אִתִּי וְנִצַּבְתָּ עַל הַצּוּר׃}
}
{
	\eJ{
		And YHWH said, “Here is a place with me, and you’ll stand
		up on a rock;
	}
}
{
%	\footnote{
%		a place with me, and you’ll stand up on a rock. Earlier, in the episode
%		of the water from the rock, God stands on the crag of Horeb, while Moses
%		stands opposite (Exod 17:6). But now Moses stands on the crag with God.
%		Whatever this means physically, it shows Moses moving closer to God. His
%		intimacy with the deity grows through the course of his life. (A student
%		of mine, Brian Kelly, observed this. Ramban saw a connection between
%		these two episodes as well.)
%	}
}

\Verse{22}
{
	\hJ{וְהָיָה בַּעֲבֹר כְּבֹדִי וְשַׂמְתִּיךָ בְּנִקְרַת הַצּוּר וְשַׂכֹּתִי כַפִּי עָלֶיךָ עַד עבְרִי׃}
}
{
	\eJ{
		and it will be, when my glory passes, that I’ll set you in
		a cleft of the rock, and I’ll cover my hand over you until
		I’ve passed,
	}
}
{
}

\Verse{23}
{
	\hJ{וַהֲסִרֹתִי אֶת כַּפִּי וְרָאִיתָ אֶת אֲחֹרָי וּפָנַי לֹא יֵרָאוּ׃}
}
{
	\eJ{
		and I’ll turn my hand away, and you’ll see my back, but my
		face won’t be seen.”
	}
}
{
}
